% ============================================================
% CHAPTER 4: SPRINT 2 --- EMPLOYEE AND LEAVE MANAGEMENT
% ============================================================

\chapter{Sprint 2: Employee and Leave Management}

\chapterplan{
    \chapterplanitem{1}{sec:s2_intro}{Introduction}
    \chapterplanitem{2}{sec:s2_analysis}{Specification and Analysis}
    \chapterplanitem{3}{sec:s2_sequence}{Sequence Diagrams}
    \chapterplanitem{4}{sec:s2_impl}{Implementation and Endpoints}
    \chapterplanitem{5}{sec:s2_ui}{Realization and Screenshots}
    \chapterplanitem{6}{sec:s2_conclusion}{Conclusion}
}

% ----------------------------------------------------------------
\section{Introduction}\label{sec:s2_intro}
% ----------------------------------------------------------------

Sprint~2 transforms the authenticated shell delivered in Sprint~1 into the first operational
Human Resources module of the platform. Once user authentication, role assignment, and
organizational reference data are available, the system can move from infrastructure concerns
to day-to-day HR operations. This sprint therefore introduces the two business capabilities that
structure the rest of the platform: employee profile management and leave request management.

The first capability establishes the employee as the central business entity of the system.
Unlike a simple user account, an employee profile connects a person to the company structure
through branch, department, designation, salary, employment start date, and administrative
documents. In the implemented platform, the employee record is linked to the underlying
\texttt{users} table and becomes the entry point reused later by payroll, insurance, social
benefits, attendance, and AI-assisted analytics.

The second capability introduces the first real approval workflow of the project. Leave requests
are not treated as static forms; they are processed through a role-aware validation chain in
which the platform distinguishes between employee submission, manager review, and, when
required, RH final approval. In code, RH roles are represented by \texttt{rh} and
\texttt{rh\_manager}; in the present chapter, these are referred to as HR roles for business
clarity. This workflow gives Sprint~2 a strategic importance: it is the sprint where the
platform begins to enforce scope, policy, and traceability rather than merely storing data.

Sprint~2 also introduces supporting services that improve operational coherence: attachment
handling for leave requests, leave balance consultation, holiday-aware working-day calculation,
attendance consultation, and in-platform notifications with unread tracking. Together, these
features turn the application into a usable HR workspace rather than a collection of isolated
CRUD pages.

% ----------------------------------------------------------------
\section{Specification and Analysis}\label{sec:s2_analysis}
% ----------------------------------------------------------------

\subsection{Sprint Goal}

The goal of Sprint~2 is to operationalize employee administration and leave processing within a
single role-governed HR workspace. The sprint is considered successful when the platform allows:
(1)~creation and maintenance of employee profiles linked to user accounts; (2)~registration and
consultation of employee-related document references; (3)~submission and review of leave
requests with policy validation; (4)~consultation of leave balances, attendance records, and
notifications; and (5)~clear separation between employee actions and management actions.

\subsection{Functional Scope}

\begin{table}[H]
\centering
\footnotesize
\renewcommand{\arraystretch}{1.25}
\begin{tabular}{|p{0.20\textwidth}|p{0.53\textwidth}|p{0.19\textwidth}|}
\hline
\thead{Workstream} & \thead{Covered Functionality} & \thead{Primary Actors} \\
\hline
\rowcolor{tablerowgray}
Employee directory & Create employee profiles, link them to existing or newly created users,
assign organizational attributes, update records, view statistics, and manage activation state
through related user actions. & Admin, RH, RH Manager \\
\hline
Employee documents & Maintain employee-document associations and uploaded document metadata
required for administrative follow-up. & RH, RH Manager \\
\hline
\rowcolor{tablerowgray}
Leave workflow & Submit leave requests, validate date ranges, prevent overlap, attach
supporting files, route approvals to manager and optionally HR, and register rejection reasons. &
Employee, Manager, RH \\
\hline
Support services & Consult leave balances, holidays, attendance records, and in-platform
notifications with unread/read state management. & All authenticated users \\
\hline
\end{tabular}
\caption{Functional Scope of Sprint 2}
\label{tab:s2_scope}
\end{table}

\subsection{Architectural Decisions and Challenges}

\begin{itemize}
    \item \textbf{Dual identity model:} the platform distinguishes between a generic
    authenticated \texttt{user} and a business \texttt{employee}. Sprint~2 must therefore
    synchronize both layers so that a person can log in while also belonging to a branch,
    department, and designation.
    \item \textbf{Employee creation as a transactional process:} employee creation does not
    simply insert one row. The backend may link an existing user, create a new one if needed,
    assign a default role, generate an \texttt{employee\_id} automatically when absent
    (\texttt{EMP-XXXXXXXX}), then persist the employee profile.
    \item \textbf{Role-aware visibility:} leave and attendance data are not exposed uniformly.
    Employees see their own requests, while managers and HR-related roles can review broader
    operational queues according to authorization logic.
    \item \textbf{Workflow modeled through business state plus approval stage:} the implemented
    leave process uses both \texttt{status} and \texttt{approval\_stage}. This is more precise
    than a single status field because it allows the system to represent an intermediate
    manager-approved state for leave types that still require HR confirmation.
    \item \textbf{Policy validation before persistence:} the leave controller validates weekends,
    company holidays, overlapping requests, leave-type maximum duration, and available balance
    before a request is inserted.
    \item \textbf{Attachment handling:} leave attachments are validated for type and size
    (\texttt{pdf}, \texttt{jpg}, \texttt{jpeg}, \texttt{png}, \texttt{webp}, max 10~MB) and
    stored through Laravel's filesystem abstraction.
    \item \textbf{Notification consistency:} each important workflow action generates an
    in-app notification stored with a type, payload, targeting data, and a
    \texttt{dedup\_key} to avoid duplicates. Read state is maintained separately through the
    \texttt{notification\_reads} table.
    \item \textbf{Frontend reactivity without WebSockets:} the Vue.js frontend relies on a
    notification store and short polling to keep unread counts and approval queues current
    without introducing real-time infrastructure in this sprint.
\end{itemize}

\subsection{Sprint Backlog}

\begin{table}[H]
\centering
\footnotesize
\renewcommand{\arraystretch}{1.25}
\begin{tabular}{|p{0.10\textwidth}|p{0.62\textwidth}|p{0.15\textwidth}|}
\hline
\thead{ID} & \thead{User Story} & \thead{Priority} \\
\hline
\rowcolor{tablerowgray}
S2-1 & As RH staff, I want to create an employee profile linked to a user account so that each
person has a complete HR identity inside the platform. & Must \\
\hline
S2-2 & As an administrator, I want to accept registered users into the employee directory and
assign their operational role during profile creation. & Must \\
\hline
\rowcolor{tablerowgray}
S2-3 & As RH staff, I want to register employee-related document references so that contracts
and administrative documents can be tracked per employee. & Must \\
\hline
S2-4 & As an employee, I want to submit a leave request with dates, duration, reason, and
optional attachments so that time-off management becomes digital. & Must \\
\hline
\rowcolor{tablerowgray}
S2-5 & As a manager, I want to review pending leave requests and perform the first validation
step according to leave type and policy. & Must \\
\hline
S2-6 & As RH staff, I want to perform final approval when a leave type requires a second level
of validation. & Must \\
\hline
\rowcolor{tablerowgray}
S2-7 & As a user, I want to consult leave balances, attendance records, and notifications so
that I can follow operational status from one workspace. & Should \\
\hline
S2-8 & As any concerned actor, I want workflow events to generate notifications with unread
tracking so that no approval or decision is missed. & Should \\
\hline
\end{tabular}
\caption{Sprint 2 Backlog}
\label{tab:sprint2_backlog}
\end{table}

\subsection{Core Use Cases}

To avoid the redundancy often produced by long, repetitive use-case descriptions, Sprint~2 is
better summarized through a compact matrix of business interactions and system responses.

\begin{table}[H]
\centering
\small
\renewcommand{\arraystretch}{1.28}
\begin{tabular}{|p{0.14\textwidth}|p{0.14\textwidth}|p{0.30\textwidth}|p{0.34\textwidth}|}
\hline
\textbf{ID} & \textbf{Actor} & \textbf{Intent} & \textbf{System Response} \\
\hline
\rowcolor{tablerowgray}
UC-S2-01 & RH / Admin & Create an employee profile from an existing or new user account. &
The system validates organizational references, generates an \texttt{employee\_id} when
missing, links or creates the \texttt{user}, assigns roles when authorized, and persists the
employee record transactionally. \\
\hline
UC-S2-02 & RH / Admin & Maintain the employee directory. &
The system lists employee records with filtering, exposes statistics, and supports activation,
suspension, or banning actions through related user-state operations. \\
\hline
\rowcolor{tablerowgray}
UC-S2-03 & Employee & Submit a leave request. &
The system validates dates, excludes weekends and holidays from working days, prevents
overlap, checks leave balance, stores attachments, and creates a pending request. \\
\hline
UC-S2-04 & Manager & Perform first-level leave review. &
The system allows approval or rejection. For ordinary leave, manager approval may complete the
workflow; for leave types such as sick leave, approval advances the request to an HR review
stage. \\
\hline
\rowcolor{tablerowgray}
UC-S2-05 & RH & Perform final leave validation when required. &
The system converts a manager-approved request into a fully approved request or records a final
rejection with explanation. \\
\hline
UC-S2-06 & Any authenticated user & Consult support information. &
The system exposes leave balances, attendance records, and notifications, and allows users to
mark notifications as read individually or in bulk. \\
\hline
\end{tabular}
\caption{Core Use Cases of Sprint 2}
\label{tab:s2_usecases}
\end{table}

\subsection{Database Design}

\begin{table}[H]
\centering
\footnotesize
\renewcommand{\arraystretch}{1.25}
\begin{tabular}{|p{0.23\textwidth}|p{0.48\textwidth}|p{0.21\textwidth}|}
\hline
\thead{Table} & \thead{Key Attributes} & \thead{Relationships / Role} \\
\hline
\rowcolor{tablerowgray}
\texttt{employees} & id, user\_id, name, email, employee\_id, branch\_id, department\_id,
designation\_id, company\_doj, salary, is\_active, created\_by & Central employee entity;
FKs to users, branches, departments, designations \\
\hline
\texttt{employee\_documents} & id, employee\_id, document\_id, file\_path, file\_name,
expiry\_date & Links employees to document definitions and stored file metadata \\
\hline
\rowcolor{tablerowgray}
\texttt{leave\_types} & id, name, description, leave\_code, maximum\_days, is\_active &
Reference catalog for leave requests and balances \\
\hline
\texttt{leaves} & id, employee\_id, leave\_type\_id, start\_date, end\_date, duration\_type,
total\_days, status, approval\_stage, reason, policy\_violations, approved\_by, rejected\_by,
rejection\_reason, approval\_probability & Core leave workflow entity; belongs to employee and
leave type \\
\hline
\rowcolor{tablerowgray}
\texttt{leave\_attachments} & id, leave\_id, file\_name, file\_path, file\_type, file\_size &
Stores metadata for uploaded leave attachments \\
\hline
\texttt{leave\_balances} & id, employee\_id, leave\_type\_id, balance
\footnotemark[1] & Validation source for remaining leave entitlement \\
\hline
\rowcolor{tablerowgray}
\texttt{holidays} & id, name, holiday\_date, description, is\_optional
\footnotemark[2] & Used to exclude company holidays from working-day computation \\
\hline
\texttt{time\_sheets} & id, employee\_id, date, check\_in, check\_out, work\_hours,
overtime\_hours, notes, status\footnotemark[3] & Operational attendance records consulted in
Sprint~2 \\
\hline
\rowcolor{tablerowgray}
\texttt{notifications} & id, type, actor\_id, target\_roles, target\_user\_ids, payload,
channel, dedup\_key & In-app workflow notifications with role and user targeting \\
\hline
\texttt{notification\_reads} & id, notification\_id, user\_id, read\_at & Tracks per-user read
state independently from notification creation \\
\hline
\end{tabular}
\caption{Sprint 2 Core Data Structures}
\label{tab:s2_db}
\end{table}

\footnotetext[1]{The \texttt{LeaveBalance} model also exposes a computed \texttt{remaining}
value used by the leave workflow.}
\footnotetext[2]{These attributes are used operationally by the leave controller and holiday
model during working-day evaluation.}
\footnotetext[3]{The attendance controller uses \texttt{status} when present in the evolved
schema.}

\subsection{Leave Approval Model}

The leave workflow is best understood as a two-level model. Final business outcome is stored in
\texttt{status} (\texttt{pending}, \texttt{approved}, \texttt{rejected}), while
\texttt{approval\_stage} records the path by which the request reached that outcome. This
design is particularly useful for leave types that require two validations.

\begin{figure}[H]
    \centering
    \imageholder{diagram}{%
        Leave approval model for Sprint 2.
        Initial state: status=pending, approval\_stage=null.
        Transition 1: Employee submits request.
        Transition 2a: Manager approves ordinary leave
            -> status=approved, approval\_stage=manager\_approved.
        Transition 2b: Manager approves sick leave
            -> status remains pending, approval\_stage=manager\_approved.
        Transition 3: RH approves previously manager-approved sick leave
            -> status=approved, approval\_stage=hr\_approved.
        Rejection path: Manager or RH rejects request
            -> status=rejected, approval\_stage=rejected, rejection reason stored.
        Terminal states: approved, rejected.
    }
    \caption{Leave Approval Model Implemented in Sprint 2}
    \label{fig:s2_fsm}
\end{figure}

This model has two advantages. First, it keeps the final state simple for reporting and search.
Second, it preserves the validation path, which is necessary when one leave type can be fully
approved by a manager while another must remain pending until RH confirmation.

\subsection{Refined Use Case Diagram}

\begin{figure}[H]
    \centering
    \imageholder{diagram}{%
        Refined use case diagram for Sprint 2.
        System boundary: Employee and Leave Management.
        Actors: Employee, Manager, RH Manager, Administrator.
        Employee use cases:
            Submit Leave Request;
            Attach Supporting File;
            Consult My Leave Requests;
            Consult Leave Balance;
            Consult Attendance;
            Read Notifications.
        Manager use cases:
            Review Pending Leave Requests;
            Approve Leave at Manager Level;
            Reject Leave Request.
        RH / Administrator use cases:
            Create Employee Profile;
            Accept Pending Registered User;
            Update Employee Information;
            Register Employee Document Metadata;
            Final Approve Leave When Required;
            Consult Attendance Statistics.
        Shared technical include:
            Authenticate via JWT;
            Apply Role / Permission Checks;
            Create In-App Notifications.
    }
    \caption{Refined Use Case Diagram of Sprint 2}
    \label{fig:s2_refined_uc}
\end{figure}

% ----------------------------------------------------------------
\section{Sequence Diagrams}\label{sec:s2_sequence}
% ----------------------------------------------------------------

\subsection{Employee Profile Creation and Directory Synchronization}

\begin{figure}[H]
    \centering
    \imageholder{diagram}{%
        Employee creation sequence diagram for Sprint 2.
        Participants: RH Manager / Admin, Vue Employee Directory, EmployeeController,
        User model, Employee model, PostgreSQL.
        Flow:
        1. Vue loads employees, branches, departments, and designations.
        2. Admin may also load users without employee profiles.
        3. Actor opens Create Employee or Accept Pending User form.
        4. Vue sends POST /api/employees with organizational data and optional user\_id.
        5. EmployeeController validates input and permissions.
        6. Transaction begins.
        7. If user\_id exists, link existing user; otherwise find or create user by email.
        8. Generate employee\_id automatically if absent.
        9. Insert employee row linked to selected branch, department, and designation.
        10. Sync user role when authorized, otherwise assign default role.
        11. Commit transaction and return created employee resource.
        12. Vue refreshes directory and shows feedback message.
    }
    \caption{Sequence Diagram: Employee Profile Creation}
    \label{fig:s2_employee_create_sequence}
\end{figure}

\subsection{Leave Submission and Approval Chain}

\begin{figure}[H]
    \centering
    \imageholder{diagram}{%
        Leave submission and approval sequence diagram for Sprint 2.
        Participants: Employee, Manager, RH, Vue Leave Requests page, LeaveController,
        LeaveBalance model, Holiday model, PostgreSQL, Notification subsystem.
        Flow:
        1. Employee opens Leave Requests page.
        2. Vue loads leave types and existing requests.
        3. Employee submits POST /api/leaves with type, duration, dates, reason,
           and optional attachments.
        4. LeaveController resolves employee scope.
        5. Working days are computed excluding weekends and holidays.
        6. Controller checks overlap, leave-type maximum duration,
           and remaining leave balance.
        7. If validation fails, return 422 with field errors.
        8. If validation succeeds, insert leave row with status=pending.
        9. Store each attachment in leave-specific storage path and create metadata rows.
        10. Create a deduplicated notification for management roles.
        11. Manager reviews the request and calls POST /api/leaves/\{id\}/approve-by-manager
            or POST /api/leaves/\{id\}/reject.
        12. For leave types requiring RH confirmation, approval\_stage becomes
            manager\_approved while status remains pending.
        13. RH performs POST /api/leaves/\{id\}/approve-by-hr when a second validation
            is required.
        14. Employee receives notification of approval or rejection.
    }
    \caption{Sequence Diagram: Leave Submission and Approval Chain}
    \label{fig:s2_leave_sequence}
\end{figure}

% ----------------------------------------------------------------
\section{Implementation and Endpoints}\label{sec:s2_impl}
% ----------------------------------------------------------------

\subsection{Implementation Approach}

Sprint~2 is implemented as a collaboration between a Laravel API backend and a Vue.js single-page
frontend. On the backend side, employee management is handled by
\texttt{EmployeeController}; leave processing is handled by \texttt{LeaveController}; attendance
consultation is handled by \texttt{AttendanceController}; and notification delivery is exposed
through \texttt{NotificationController}. The implementation combines classical CRUD operations
with custom workflow endpoints for approval actions.

On the frontend side, the employee directory and leave requests pages act as operational
workspaces rather than simple forms. The employee directory integrates record creation, pending
user acceptance, status management, and statistical consultation. The leave page combines
submission, queue consultation, approval actions, and attachment handling in the same interface.
Notifications are centralized in a Pinia store and refreshed periodically, which keeps the
approval cycle usable without a WebSocket infrastructure.

\subsection{API Endpoint Specification}

\begin{table}[H]
\centering
\small
\renewcommand{\arraystretch}{1.26}
\begin{tabular}{|p{0.08\textwidth}|p{0.34\textwidth}|p{0.44\textwidth}|}
\hline
\thead{Method} & \thead{Endpoint URI} & \thead{Purpose} \\
\hline
\rowcolor{tablerowgray}
GET & \texttt{/api/employees} & List employee profiles with optional search, branch,
department, and designation filters. \\
\hline
POST & \texttt{/api/employees} & Create an employee profile and link or create the underlying
user account. \\
\hline
\rowcolor{tablerowgray}
GET & \texttt{/api/employees/\{id\}} & Retrieve one employee profile with organizational
relations. \\
\hline
PUT / PATCH & \texttt{/api/employees/\{id\}} & Update employee information. \\
\hline
\rowcolor{tablerowgray}
DELETE & \texttt{/api/employees/\{id\}} & Remove an employee record when authorized. \\
\hline
GET & \texttt{/api/employees/\{id\}/statistics} & Return employee-level statistics used by
the RH directory and analytics cards. \\
\hline
\rowcolor{tablerowgray}
GET / POST & \texttt{/api/employees/employee-documents} & List or create employee-document
association records. \\
\hline
GET / POST & \texttt{/api/employees/document-uploads} & Manage generic uploaded document
metadata referenced by employee workflows. \\
\hline
\rowcolor{tablerowgray}
GET & \texttt{/api/leaves} & List leave requests with role-aware visibility. \\
\hline
POST & \texttt{/api/leaves} & Submit a leave request with policy validation and optional
attachments. \\
\hline
\rowcolor{tablerowgray}
GET & \texttt{/api/leaves/\{id\}} & Retrieve one leave request. \\
\hline
PUT / PATCH & \texttt{/api/leaves/\{id\}} & Update an existing leave request when permitted. \\
\hline
\rowcolor{tablerowgray}
DELETE & \texttt{/api/leaves/\{id\}} & Delete a leave request when permitted. \\
\hline
POST & \texttt{/api/leaves/\{id\}/attachments} & Upload one or more files associated with a
leave request. \\
\hline
\rowcolor{tablerowgray}
GET & \texttt{/api/leaves/\{leaveId\}/attachments/\{attachmentId\}/download} &
Download a stored leave attachment. \\
\hline
POST & \texttt{/api/leaves/\{id\}/approve-by-manager} & Perform the manager approval step. \\
\hline
\rowcolor{tablerowgray}
POST & \texttt{/api/leaves/\{id\}/approve-by-hr} & Perform final RH approval when the leave
type requires it. \\
\hline
POST & \texttt{/api/leaves/\{id\}/reject} & Reject a leave request and persist the reason. \\
\hline
\rowcolor{tablerowgray}
GET / POST & \texttt{/api/leaves/types} & Consult or manage the leave type catalog. \\
\hline
GET / POST & \texttt{/api/leaves/balances} & Consult or manage leave balances. \\
\hline
\rowcolor{tablerowgray}
GET / POST & \texttt{/api/leaves/holidays} & Consult or manage company holidays used in
working-day evaluation. \\
\hline
GET & \texttt{/api/attendance} & List attendance data in a role-scoped manner. \\
\hline
\rowcolor{tablerowgray}
GET / POST & \texttt{/api/attendance/records} & Consult or create attendance records. \\
\hline
PUT / DELETE & \texttt{/api/attendance/records/\{id\}} & Update or delete attendance records. \\
\hline
\rowcolor{tablerowgray}
GET & \texttt{/api/attendance/statistics} & Return aggregated attendance indicators for a
selected period. \\
\hline
GET & \texttt{/api/notifications} & List visible in-app notifications for the authenticated
user. \\
\hline
\rowcolor{tablerowgray}
POST & \texttt{/api/notifications/\{id\}/mark-read} & Mark one notification as read. \\
\hline
POST & \texttt{/api/notifications/mark-all-read} & Mark all visible notifications as read. \\
\hline
\rowcolor{tablerowgray}
GET & \texttt{/api/notifications/unread-count} & Return the number of unread notifications. \\
\hline
\end{tabular}
\caption{Sprint 2 Core Endpoint Specification}
\label{tab:s2_endpoints}
\end{table}

\paragraph{Complementary implementation note.}
The current platform also exposes \texttt{/api/leaves/request-insights} and
\texttt{/api/leaves/optimal-dates} as non-blocking AI-assisted helpers. These enrich the user
experience of Sprint~2 but remain complementary to the core business scope described in this
chapter.

% ----------------------------------------------------------------
\section{Realization and Screenshots}\label{sec:s2_ui}
% ----------------------------------------------------------------

\begin{figure}[H]
    \centering
    \imageholder{screenshot}{%
        Employee Directory screen (Sprint 2, RH/Admin workspace).
        Header: icon, title "Employee Directory", subtitle, and "Add Employee" button.
        Optional alert banner for registered users waiting for a full employee profile.
        Creation form fields: name, email, gender, role, branch, department, designation,
        salary, and address. Main table columns: employee ID, name, email, salary, status.
        Row actions: AI Risk, Suspend, Ban, Activate, or Accept for pending profiles.
        Right-side insight card shows turnover prediction and recommendations for the selected
        employee.
    }
    \caption{Sprint 2 --- Employee Directory and Profile Creation}
    \label{fig:s2_employee_ui}
\end{figure}

\begin{figure}[H]
    \centering
    \imageholder{screenshot}{%
        Leave Requests screen (Sprint 2, shared employee and manager workspace).
        Top hero section shows request counters and current access mode.
        Left card: leave submission form with employee selector for manager view, leave type,
        duration toggle, start date, end date, selected range summary, reason textarea,
        and multi-file attachment input.
        Right card: searchable queue overview table with employee, period, duration, status,
        and reason. Supporting panels explain approval workflow and display AI suggestions.
        For management roles, a lower approval area exposes Approve and Reject buttons for
        pending requests.
    }
    \caption{Sprint 2 --- Leave Submission and Approval Workspace}
    \label{fig:s2_leave_ui}
\end{figure}

\begin{figure}[H]
    \centering
    \imageholder{screenshot}{%
        Notifications screen (Sprint 2).
        Title "Notifications", refresh action, "Mark All as Read" button,
        unread counter, and filter chips for All and Unread.
        Main list displays notification cards with title, message, timestamp, and action link.
        Read status is individualized per user and synchronized with the top-bar notification bell.
    }
    \caption{Sprint 2 --- Notification Center}
    \label{fig:s2_notifications_ui}
\end{figure}

% ----------------------------------------------------------------
\section{Conclusion}\label{sec:s2_conclusion}
% ----------------------------------------------------------------

Sprint~2 marks the transition from a technically prepared platform to an operational HR system.
Its main contribution is not only the addition of new tables and pages, but the installation of a
reusable pattern for the rest of the project: a central employee entity, role-scoped business
actions, explicit validation rules, and in-app feedback through notifications.

The employee directory created in this sprint gives the application its first real business
backbone. Every major functional area that follows---payroll, insurance, benefits, and
analytics---relies on the existence of a coherent employee record linked to the authentication
layer and the organizational structure. The leave module, in turn, introduces approval-chain
logic, policy-aware validation, attachment support, and operational traceability.

From an architectural perspective, Sprint~2 is also where the platform begins to combine backend
authorization with frontend usability. The same Vue workspaces support both data entry and queue
supervision, while Laravel enforces scope and workflow rules at API level. This separation of
concerns is essential for the next iterations of the project.

Sprint~3 can therefore build on a stable business foundation: employee identities are
established, leave workflows are controlled, attendance consultation is available, and
notifications already connect actions to users. The platform is now ready to extend its HR
coverage toward payroll production, social-service workflows, and more advanced decision support.
